\section{Discussion}\label{sec:discussion}

\paragraph{What a certificate is and is not.}
A certified kernel is a statement about a declared class: a registered map, a floor class, a bounded defect and a horizon. It is not a prediction of the trajectory and it is not a probability. The stochastic readings of Sections~\ref{sec:cod-stoch} and~\ref{sec:edw-stoch} are reported beside the kernels so that a reader can see what the two kinds of statement cost: the i.i.d.\ bootstrap survival of the cod LRP under BAU is $0.90$ over twenty years with a ceiling of $0.906$ set by the two collapse residuals, and the frozen-floor bootstrap of $\Cstar$ has $12.5\%$ of its mass at or below zero; the attractor bootstrap puts BAU at the 618-ft trigger at $0.41$. None of these numbers is the kernel and none contradicts it. The kernel says what holds for every sequence in the class; the bootstrap says how often it holds for the sequences the record generates.

\paragraph{The declared defect.}
One audit exceeds its declaration out of sample and the other does not. The cod out-of-sample residuals (2008--2015) are small, with maximum $46.7\kt$, and the declaration of $329\kt$ is conservative by a factor of seven for the post-moratorium regime; the Edwards out-of-sample maximum of $21.8\ft$ exceeds the declared $15.4\ft$; the positive-pumping horizon falls from three years to two, while zero pumping retains three. We keep the declarations as registered and report the audited horizons beside them, because the alternative---redeclaring the defect after seeing the data it is meant to bound---is the practice the protocol exists to prevent. Table~\ref{tab:edw-unc} shows what redeclaration would cost: the supremum defect the three-year aquifer kernel can carry is $16.0$--$17.6\ft$ for the pumping rules ($24.0\ft$ at zero pumping), so a declaration at the out-of-sample maximum leaves only zero pumping certified through $T=3$; even zero pumping fails at $T=4$.

\paragraph{Where the cod verdict is fragile.}
The cod negatives are fragile in one direction and robust in the other. In the fragile direction, the vacuous class under the worst floor ($\Cvac<0$: no state is one-step safe at the 1992 residual) disappears once a share $\lambda\ge0.2$ of the residual variance is attributed to observation error, under a window that starts after the moratorium ($\Cstar=165.7\kt$ for 1995--2015), under the Fox form ($\Cstar=79.1\kt$) and under the 2024 assessment's lower reference point ($\Cstar=-7.7\kt$ at 2005--2024, $+8.0\kt$ at 1995--2024). In the robust direction, the no-dominance verdict for the Stage-I rule does not depend on $K$ (Section~\ref{sec:cod-ident}), on the floor, on the form or on the window, because it is a comparison of two rules that pump the same catch at the LRP. We regard the second kind of result as the paper's contribution to the cod case and the first kind as its honest scope.

\paragraph{Where the aquifer verdict is fragile.}
The Edwards certificate at 618 ft is a statement about a three-parameter affine map with a worst recharge floor taken from the fit window. The map omits spring discharge, the storage nonlinearity of the karst and the spatial structure of pumping; its out-of-sample residual is larger than its fit residual; and its 1950s replay reproduces the drought minimum with an $8\ft$ bias in the unsafe direction. The attractor ordering of the rules survives all of this because it depends only on the sign of $\gamma$, which is negative in every one of 5000 bootstrap refits (the flat-90\% minus BAU gap is positive in all of them, with minimum $0.5\ft$); the location of the BAU attractor relative to 618 ft, a difference of $2.3\ft$ against a bootstrap standard deviation of $13\ft$, does not, and the unsecurability of 660 ft, a shortfall of $2.1$--$12.7\ft$ against a standard deviation of $10\ft$ on the zero-pumping attractor, is a statement about the registered fit rather than about the aquifer. A reader who wants a finer resolution should fit a finer map; the frame will certify it on the same terms.

\paragraph{The retention rule as a design instrument.}
Clause (b) of the retention rule is a supply comparison at equal protection, and it is the only place in the frame where a staged rule can earn its place. The two systems show the two outcomes: where the observed states are mostly above the stage boundaries (the aquifer, 65\% of years above 660 ft), a staged rule banks supply in the unconstrained years and is retained; where they are mostly below (cod, 83\% of the fit-window years below the LRP, where the staged rules take little or nothing), the staged rule supplies less than the flat cap it matches in protection and is not. A manager designing a staged rule can read the rule's prospects off the empirical distribution of the state relative to its boundaries before any kernel is computed.

\paragraph{Limits of the one-dimensional setting.}
Everything above uses that the state is a scalar and the safe set an interval, which makes the kernels half-lines and the exact bounded-defect recursion a one-dimensional dynamic programme. The frame extends to monotone maps on ordered state spaces without change in the statements and with the kernels becoming upper sets; the exact recursion becomes a robust dynamic programme on a grid and loses the closed forms. The staged-rule proposition and the securability theorem are the two results that depend most on dimension one, through the fixed-point structure of Theorem~\ref{thm:kernel}; their analogues in higher dimension are conjectures we do not make here.
